\section{The \uvif{} Organizational Framework}
\label{sec:uvif_framework}

\subsection{State, fields, and admissible dynamics}
Let $\mathbf{s}\in\mathcal{S}\subseteq\mathbb{R}^n$ describe an explicitly chosen operating state. Depending on the application, its coordinates may include belief summaries, policy parameters, protected memory variables, and resource reserves. The state must contain the variables needed to model their evolution; otherwise the description may require additional memory or stochastic inputs. The following deterministic formulation is a modelling template, not a complete agent implementation.

Write the unconstrained update field as
\begin{equation}
 F(\mathbf{s},t)=w_E\Epull-w_D\Drag-w_P\Prot
                   +w_A\Eff+w_S\Sus,
 \label{eq:uvif_dynamics}
\end{equation}
where each field is evaluated at $(\mathbf{s},t)$ and each weight $w_i=w_i(\mathbf{s},t)\geq0$ is specified by the model. All weighted terms must be vectors in the same state coordinates and have units of state change per unit time. Raw information gain, latency, and energy consumption cannot be added without explicit scaling and a mapping into these coordinates. The signs are conventions: $\Drag$ and $\Prot$ are defined in the direction of increasing cost or damage, so their negative contributions oppose those changes.

State-dependent weights may encode resource depletion, urgency, or threat. Their update rules are part of the model and must be reported or estimated; dependence on state does not make them free of modelling assumptions. A fixed-weight model is a legitimate special case. A component may also vanish in a particular context, and balance does not require equal force magnitudes.

For a nonempty, closed, convex, time-independent admissibility set $\Adm$, one candidate constrained dynamics is
\begin{equation}
 \dot{\mathbf{s}}=G(\mathbf{s},t)
 :=\Pi_{T_{\Adm}(\mathbf{s})} F(\mathbf{s},t),
 \qquad \mathbf{s}(0)\in\Adm,
 \label{eq:projected}
\end{equation}
where $T_{\Adm}(\mathbf{s})$ is the tangent cone of feasible directions and $\Pi$ is Euclidean projection onto that cone. This is a first-order idealization in which the update velocity can be selected directly. Existence, uniqueness, and invariance require appropriate regularity of the field and a suitable solution concept for the projected dynamics. It does not automatically implement a feasible controller for a plant with actuator limits, delays, inertia, or uncertain dynamics.

\begin{figure}[tbp]
\centering
\renewcommand{\baselinestretch}{1}\small
\begin{tikzpicture}[box/.style={draw,rounded corners,align=center,text width=4.2cm,minimum height=1.1cm,font=\small},arr/.style={-{Latex},thick}]
\node[box] (e) at (-2.7,4) {Epistemic pull\\Learning};
\node[box] (p) at (2.7,4) {Selective protection\\Retention};
\node[box] (d) at (-2.7,2) {Computational drag\\Inference cost};
\node[box] (a) at (2.7,2) {Action efficiency\\Timely action};
\node[box] (s) at (0,0) {Sustainability\\Future capacity};
\node[box,text width=10cm] (f) at (0,-1.7) {Compatible weighted update field $F(\mathbf{s},t)$};
\node[box,text width=10cm] (g) at (0,-3.2) {Admissibility mechanism and feasible dynamics $G$};
\draw[arr] (e.west) -- ++(-0.45,0) |- (f.west);
\draw[arr] (d.south) -- (f.north -| d.south);
\draw[arr] (p.east) -- ++(0.45,0) |- (f.east);
\draw[arr] (a.south) -- (f.north -| a.south);
\draw[arr] (s.south) -- (f.north);
\draw[arr] (f) -- (g);
\end{tikzpicture}
\caption{Five-component modelling architecture. Hard admissibility is distinct from soft protection and cost terms. Arrows indicate model composition, not a proven causal decomposition. Schematic prepared as editable TikZ code with OpenAI Codex (GPT-6) assistance.}
\label{fig:uvif_force_architecture}
\end{figure}


\subsection{The five regulatory components}
\label{subsec:epistemic}
\textbf{Epistemic pull.} This component favours informative changes in the operating state. In a differentiable surrogate model one may set
\begin{equation}
 \Epull(\mathbf{s},t)=\nabla_{\mathbf{s}} I(\mathbf{s},t),
 \label{eq:epistemic_pull}
\end{equation}
where $I$ is expected information gain under a stated observation model and sampling policy. Expected information gain must identify the uncertainty being reduced; it is not interchangeable with confidence or task reward. For discrete decisions, a gradient requires a justified relaxation or an alternative update rule. Curiosity and active inference provide motivation for this component \citep{dubey2019curiosity,parr2018generalised}.

\textbf{Selective protection.} A candidate field is $\Prot=\nabla_{\mathbf{s}}L_P$, where $L_P$ measures damage to specified internal structures or loss of retained competence. This can describe regularization or resistance to updating protected parameters. It may reduce interference with previous learning, but it does not guarantee freedom from catastrophic forgetting. Protection also differs from a hard safety requirement: a soft penalty can be outweighed, whereas an admissibility constraint is intended to exclude a state or action.

\textbf{Computational drag.} A compatible cost-based field is
\begin{equation}
 \Drag(\mathbf{s},t)=\nabla_{\mathbf{s}} C_{\mathrm{comp}}(\mathbf{s},t),
 \label{eq:drag}
\end{equation}
where $C_{\mathrm{comp}}$ is a differentiable surrogate for computation, latency, or memory expenditure. This is a cost gradient, not literal mechanical dissipation. A velocity-dependent damping term would require a different, explicitly specified dynamical model. Resource-rational analysis supplies an established basis for making computational expenditure part of decision-making \citep{lieder2019resourcerational}.

\textbf{Action efficiency.} This component rewards useful and timely action. For a differentiable policy $\pi_{\mathbf{s}}$ and transition model $p$, define
\begin{align}
 J_A(\mathbf{s},t)
 &=\mathbb{E}_{a\sim\pi_{\mathbf{s}},\,\mathbf{s}'\sim p(\cdot\mid\mathbf{s},a)}
       [r(\mathbf{s},a,\mathbf{s}')-c_{\mathrm{act}}(\mathbf{s},a)],
       \label{eq:action_value}\\
 \Eff(\mathbf{s},t)&=\nabla_{\mathbf{s}}J_A(\mathbf{s},t).
 \label{eq:action_eff}
\end{align}
Here $r$ is an independently specified task benefit and $c_{\mathrm{act}}$ represents action cost or delay. Defining value only as proximity to an as-yet unknown equilibrium would be circular. This example is a local surrogate; a sequential implementation must specify its planning horizon and account for state transitions. Costs assigned to this component should not also be counted in computational drag without an explicit reason.

\textbf{Sustainability.} This component represents prospective loss of operating capacity, for example through
\begin{equation}
 \Sus(\mathbf{s},t)=-\nabla_{\mathbf{s}}L_S(\mathbf{s},t),
 \label{eq:sustainability}
\end{equation}
where $L_S$ evaluates predicted depletion or degradation over a stated horizon. A penalty on instantaneous resource use alone does not establish long-term sustainability. A reserve $b$, for example, may need a balance law
\begin{equation}
 \dot b=r_{\mathrm{in}}(\mathbf{s},a,t)-r_{\mathrm{use}}(\mathbf{s},a,t),
 \qquad b\geq b_{\min},
 \label{eq:reserve}
\end{equation}
with replenishment and consumption rates in compatible units. If $b$ is part of $\mathbf{s}$, this balance law must be implemented in the corresponding coordinate of the full dynamics; it cannot be an unrelated auxiliary equation. Projection alone cannot create an unavailable replenishment action.

The distinction between drag and sustainability is temporal and operational: drag concerns present inference expenditure, whereas sustainability concerns future availability of required resources. Selective protection concerns damage to designated structures. These quantities can overlap, so the decomposition requires an accounting convention, normalization, and evidence that the components can be separately measured or manipulated.

\subsection{Admissibility and safety enforcement}
\label{subsec:admissibility}
\begin{definition}[Admissibility set]
$\Adm$ is the set of states satisfying the physical and operational requirements chosen for a specified deployment context. Where smooth inequality constraints suffice, write $\Adm=\{\mathbf{s}:h_j(\mathbf{s})\geq0,\ j=1,\ldots,m\}$.
\end{definition}

At a smooth boundary, an inward or tangent velocity satisfies $\nabla h_j(\mathbf{s})^\top\dot{\mathbf{s}}\geq0$ for active constraints. Such conditions motivate projection and barrier-based safety mechanisms, but their applicability depends on the actual system dynamics \citep{ames2019barrier}. A time-dependent constraint additionally involves $\partial h_j/\partial t$. Nonconvex domains, discrete actions, and uncertain disturbances require separate treatment.

Ethical and legal requirements can inform admissibility specifications, but they are not generally reducible to a complete set of state inequalities. A mathematical set certifies only the encoded model properties under its assumptions. It does not establish ethical acceptability or compliance by itself. Risk is handled through protection, explicit risk constraints, or the task model; it is not introduced as an unspecified sixth force.

\subsection{Equilibrium, stability, and viable operation}
\begin{definition}[Stationary Variational Equilibrium]
For an autonomous constrained dynamics $\dot{\mathbf{s}}=G(\mathbf{s})$, the equilibrium set is
\begin{equation}
 \VEq=\{\mathbf{s}^*\in\Adm:G(\mathbf{s}^*)=0\}.
 \label{eq:veq}
\end{equation}
A stationary Variational Equilibrium is any element of this set. Stability is an additional property, not a consequence of this definition.
\end{definition}

An equilibrium is Lyapunov stable if, for every $\varepsilon>0$, a sufficiently small initial displacement keeps the trajectory within $\varepsilon$ for all future times. It is locally asymptotically stable if it is stable and trajectories from a neighbourhood also converge to it. Initial states and neighbourhoods are understood relative to $\Adm$. For a nonisolated equilibrium set, convergence to the set need not imply convergence to a particular equilibrium.

At an interior equilibrium, $G=F=0$. At a constrained boundary equilibrium, $G=0$ can hold even when $F\neq0$, because the unconstrained update points out of the feasible region. The equilibrium set need not be smooth, connected, nonempty, or a manifold. Moreover, a trajectory cannot move along a set of stationary points under the same autonomous field while $G$ remains zero everywhere on that trajectory.

An invariant operating set $K\subseteq\Adm$ is a different object: trajectories starting in $K$ remain there, potentially with $G\neq0$. It can contain cycles or other sustained activity. Viability asks whether admissible operation can be maintained by available controls, whereas invariance is relative to specified dynamics or a controller. Neither property is equivalent to zero net force or Pareto optimality.

For time-varying environments, the instantaneous zero set $\VEq(t)$ need not itself be invariant. Tracking a selected $\mathbf{s}^*(t)$ requires a controller and bounds on disturbances and the rate of change of that reference. Accordingly, the broader \uvif{} design objective is adaptive viability, and stationary equilibrium is one restricted way to study it.

\subsection{A restricted analytical example}
\label{subsec:example}
Consider a dimensionless scalar operating variable $x\in[0,1]$ in a fixed environment. Let $e,p,a,s\in[0,1]$ be stipulated reference levels and let $\alpha_E,\alpha_P,\alpha_D,\alpha_A,\alpha_S>0$ be rate coefficients. Define the weighted component fields by
\begin{align}
 w_E\Epull&=\alpha_E(e-x), & w_P\Prot&=\alpha_P(x-p),\nonumber\\
 w_D\Drag&=\alpha_Dx, & w_A\Eff&=\alpha_A(a-x),\nonumber\\
 w_S\Sus&=\alpha_S(s-x).&&
 \label{eq:toy_fields}
\end{align}
These linear fields are illustrative choices, not measured descriptions of cognition. Writing $k=\alpha_E+\alpha_P+\alpha_D+\alpha_A+\alpha_S$, their net field is $F(x)=k(x^*-x)$, where
\begin{equation}
 x^*=\frac{\alpha_E e+\alpha_P p+\alpha_A a+\alpha_S s}{k}.
 \label{eq:toy_equilibrium}
\end{equation}
The stated assumptions imply $x^*\in[0,1]$. The field points inward or is tangent at both endpoints, so no projection correction is needed. For any $x(0)\in[0,1]$,
\begin{equation}
 x(t)=x^*+[x(0)-x^*]e^{-kt},
 \label{eq:toy_solution}
\end{equation}
which remains in the interval and converges exponentially to $x^*$. Equivalently, $V=(x-x^*)^2/2$ satisfies $\dot V=-k(x-x^*)^2$. This proves stability only for this stipulated scalar example.

The same dynamics is gradient descent on
\begin{align}
 \Phi(x)=\tfrac12[&\alpha_E(x-e)^2+\alpha_P(x-p)^2+\alpha_Dx^2\nonumber\\
                 &+\alpha_A(x-a)^2+\alpha_S(x-s)^2].
 \label{eq:toy_potential}
\end{align}
Thus its equilibrium also minimizes a scalar objective on $[0,1]$. The example demonstrates that a force decomposition can be compatible with optimization; it cannot establish superiority over it. More generally, fixed weights and compatible gradient fields can yield a common potential, whereas state-dependent weights need not. Determining the relevant case is a property of a specified implementation.
